\chapter{Resultados y Discusión}

\section{Tabla consolidada de resultados}

Los doce modelos evaluados —seis escenarios $\times$ dos métodos de
imputación— producen los resultados recogidos en la
Tabla~\ref{tab:resultados_mcar}. Cada fila corresponde a un nivel de
incompletitud distinto; los valores reportados son Accuracy y F1-score
medidos sobre las 655 instancias del conjunto de prueba con el clasificador
de Regresión Logística (\texttt{class\_weight=‘balanced'}, \texttt{SEED = 42}).

\begin{table}[H]
\centering
\caption{Accuracy y F1-score por método de imputación y porcentaje de
         datos faltantes MCAR. Experimento sobre \textit{Water Potability Dataset}
         con Regresión Logística (\texttt{SEED = 42}).}
\label{tab:resultados_mcar}
\begin{tabular}{@{}ccccc@{}}
\toprule
\multirow{2}{*}{\textbf{Faltantes (\%)}}
  & \multicolumn{2}{c}{\textbf{Accuracy}}
  & \multicolumn{2}{c}{\textbf{F1-score}} \\
\cmidrule(lr){2-3} \cmidrule(lr){4-5}
  & \textbf{Media} & \textbf{KNN} & \textbf{Media} & \textbf{KNN} \\
\midrule
 5  & 0.5351 & \textbf{0.5488} & 0.4857 & \textbf{0.4897} \\
10  & 0.5305 & \textbf{0.5335} & 0.4690 & \textbf{0.4814} \\
20  & \textbf{0.5107} & 0.5076 & \textbf{0.4659} & 0.4553 \\
30  & \textbf{0.5412} & 0.5213 & \textbf{0.4872} & 0.4767 \\
40  & 0.5244 & \textbf{0.5366} & \textbf{0.4800} & 0.4704 \\
50  & \textbf{0.5137} & 0.5015 & \textbf{0.4491} & 0.4333 \\
\bottomrule
\multicolumn{5}{l}{\small Tabla 4.1: Comparación de Accuracy y F1-score entre la
imputación por media y la imputación KNN, según el porcentaje de datos
faltantes introducidos bajo el mecanismo MCAR.}
\end{tabular}
\end{table}

\section{Análisis por métrica}

\subsection{Accuracy}

La Figura~\ref{fig:accuracy_mcar} traza la evolución del Accuracy de
ambos métodos a lo largo de los seis escenarios. Con 5\% y 10\% de datos
faltantes, KNN obtiene Accuracy superior al de la imputación por media
($+$0.0137 y $+$0.0030, respectivamente). A partir del 20\%, esa relación
se invierte: la imputación por media iguala o supera a KNN en cuatro de
los cuatro escenarios restantes.

La brecha más amplia a favor de la media se registra al 30\% MCAR
($\Delta = -0.0199$ para KNN), mientras que la mayor ventaja de KNN ocurre
en el escenario de menor incompletitud ($\Delta = +0.0137$ al 5\%). En
cualquier caso, ambos métodos se mantienen dentro del rango 0.50--0.55,
valor que refleja el techo de rendimiento impuesto por la baja
separabilidad lineal del \textit{dataset}.

\begin{figure}[H]
\centering
\includegraphics[width=0.95\textwidth]{chapters/img/mcar_accuracy_f1.png}
\caption{Accuracy (izquierda) y F1-score (derecha) en función del
         porcentaje de datos faltantes MCAR para los dos métodos de
         imputación evaluados.}
\label{fig:accuracy_mcar}
\end{figure}

\subsection{F1-score}

El F1-score replica el patrón anterior y lo acentúa en los extremos.
KNN aventaja a la media en los escenarios del 5\% ($+$0.0040) y 10\%
($+$0.0124), pero desde el 20\% en adelante la media retoma el liderazgo
en todos los escenarios sin excepción.

La divergencia más notable ocurre al 50\% de datos faltantes: la media
alcanza un F1 de 0.4491 frente a 0.4333 de KNN, una diferencia de 0.0158.
Cuando la mitad de los valores han sido eliminados, el espacio de atributos
disponibles para construir vecindarios resulta demasiado reducido, y KNN
pierde la ventaja que la estructura local del \textit{dataset} podría
haberle otorgado en condiciones de mayor completitud.

\begin{figure}[!htbp]
\centering
\includegraphics[width=0.88\textwidth]{chapters/img/mcar_diferencia_metodos.png}
\caption{Diferencia en Accuracy (izquierda) y F1-score (derecha) entre
         KNN e imputación por media ($\Delta = \text{KNN} - \text{Media}$)
         por porcentaje de datos faltantes. Barras en rojo indican
         ventaja de KNN; barras en azul indican ventaja de la media.}
\label{fig:diferencia_mcar}
\end{figure}

\section{Discusión}

\subsection{Efecto del porcentaje de datos faltantes sobre el rendimiento}

La hipótesis H1 postulaba una degradación progresiva del rendimiento
conforme aumenta el porcentaje de datos faltantes. Los resultados no
respaldan esa tendencia: tanto el Accuracy como el F1-score oscilan de
forma irregular entre el 5\% y el 50\% de incompletitud, sin una
caída sostenida para ninguno de los dos métodos.

Este comportamiento tiene una explicación teórica directa. Bajo el
mecanismo MCAR, la media muestral es un estimador insesgado del valor
faltante, por lo que el clasificador no acumula sesgo sistemático a medida
que crece la proporción de imputaciones \parencite{little2019}. La
varianza observada entre escenarios responde más a la aleatoriedad de
la máscara MCAR —evaluada con una única semilla— que a un efecto real
del porcentaje de incompletitud sobre el rendimiento del modelo. A ello
se suma que ninguna correlación entre las variables fisicoquímicas del
\textit{dataset} y la etiqueta de potabilidad supera 0.15 en valor
absoluto, lo que establece un techo de rendimiento bajo para la
Regresión Logística independientemente de cómo se imputen los valores
faltantes.

\subsection{Comparación entre métodos por porcentaje de faltantes}

La hipótesis H2 recibe respaldo parcial. En los escenarios de baja
incompletitud (5\% y 10\%), KNN supera a la imputación por media tanto
en Accuracy como en F1-score, lo que concuerda con la lógica de los
métodos basados en similitud: cuando la mayor parte de los atributos
está disponible, los vecindarios son representativos y la estimación
local supera a la estimación global \parencite{troyanskaya2001,garcia_laencina2010}.

Desde el 20\% de datos faltantes, sin embargo, la imputación por media
se vuelve consistentemente más robusta en F1-score. \textcite{vanbuuren2018}
documenta precisamente esta dinámica: bajo MCAR con clasificadores lineales,
la media es el estimador óptimo en esperanza, y su robustez relativa
crece a medida que la reducción del espacio de atributos observados
perjudica la calidad de los vecindarios que construye KNN.

\subsection{Implicancias prácticas}

Los resultados apuntan a tres conclusiones de uso directo en el diseño de
\textit{pipelines} con datos incompletos:

\begin{enumerate}
    \item \textbf{Datos escasamente incompletos (5\%--10\%):} KNN puede
          rendir ligeramente mejor que la imputación por media, y ese
          margen puede justificar el mayor costo computacional en
          aplicaciones donde cada décima de F1 importa.

    \item \textbf{Datos con incompletitud moderada o alta (20\%--50\%):}
          la imputación por media ofrece resultados equivalentes o
          superiores con una fracción del tiempo de cómputo. Bajo MCAR
          y con un clasificador lineal, añadir complejidad al paso de
          imputación no se traduce en mejora del rendimiento final.

    \item \textbf{El cuello de botella no es la imputación:} la diferencia
          de F1-score entre el mejor escenario (0.4897 con KNN al 5\%) y
          el peor (0.4333 con KNN al 50\%) apenas alcanza 0.056. El límite
          real de rendimiento está impuesto por la separabilidad del
          \textit{dataset}, no por el método de imputación elegido.
\end{enumerate}

\subsection{Contraste con la hipótesis original}

La pregunta inicial de la investigación —si KNN produce mejor rendimiento
que la imputación por media— no obtiene una respuesta afirmativa general.
KNN gana en los dos escenarios de menor incompletitud, pierde en tres de
los cuatro restantes y empata con la media en ninguno. Las diferencias,
además, son reducidas: en la mayoría de escenarios no superan 0.02 puntos
ni en Accuracy ni en F1-score.

Esto no descalifica a KNN como método de imputación; lo que sí queda en
evidencia es que su ventaja sobre la media depende del porcentaje de
datos ausentes, de la estructura de correlación del \textit{dataset} y
del tipo de clasificador empleado. Al replantear la pregunta de
investigación hacia el efecto del porcentaje de incompletitud, el trabajo
captura esa condicionalidad, algo que una comparación de punto único al
10\% habría pasado por alto.

\section{Limitaciones y trabajo futuro}

Cuatro aspectos limitan el alcance de las conclusiones y abren líneas de
continuación directas:

\begin{itemize}
    \item Los experimentos emplean una única semilla (\texttt{SEED = 42}).
          Repetir el diseño con varias semillas y calcular intervalos de
          confianza por \textit{bootstrap} permitiría determinar si las
          diferencias observadas entre métodos son estadísticamente
          significativas o atribuibles a la variabilidad del muestreo
          aleatorio.

    \item El uso exclusivo de Regresión Logística restringe la
          generalización. Clasificadores no lineales como Random Forest
          o XGBoost pueden explotar mejor la información recuperada por
          KNN, especialmente a bajos porcentajes de incompletitud, y
          podrían cambiar las conclusiones sobre qué método es preferible.

    \item El análisis cubre únicamente el mecanismo MCAR. Extenderlo a
          MAR y MNAR —mecanismos más frecuentes en datos reales— ampliaría
          tanto la validez externa del experimento como su relevancia
          práctica.

    \item Incorporar métodos de imputación adicionales (regresión lineal,
          MICE, \textit{missForest}) daría más contexto a los resultados
          y ubicaría a KNN y a la media dentro de un espectro más
          representativo de las opciones disponibles.
\end{itemize}
